\documentclass[11pt,a4paper]{article}

% ============================================================
% DS ANALYSE PCSI -- SUITES, CONTINUITE, DERIVATION
% Niveau solide PCSI / début CCP-INP
% ============================================================

\usepackage[utf8]{inputenc}
\usepackage[T1]{fontenc}
\usepackage[french]{babel}
\usepackage{lmodern}
\usepackage{microtype}
\usepackage{geometry}
\geometry{margin=1.7cm}

\usepackage{amsmath,amssymb,amsthm}
\usepackage{mathtools}
\usepackage{array}
\usepackage{tabularx}
\usepackage{booktabs}
\usepackage{enumitem}
\usepackage{xcolor}
\usepackage{fancyhdr}
\usepackage{lastpage}
\usepackage{titlesec}
\usepackage{tcolorbox}
\usepackage{hyperref}

\hypersetup{
    colorlinks=true,
    linkcolor=blue!60!black,
    urlcolor=blue!60!black
}

% ============================================================
% COULEURS
% ============================================================

\definecolor{mainblue}{RGB}{30,70,130}
\definecolor{darkred}{RGB}{140,30,30}
\definecolor{darkgreen}{RGB}{20,110,70}
\definecolor{lightblue}{RGB}{235,242,252}
\definecolor{lightred}{RGB}{252,238,238}
\definecolor{lightgreen}{RGB}{238,250,244}
\definecolor{lightgray}{RGB}{245,245,245}

% ============================================================
% EN-TETES
% ============================================================

\pagestyle{fancy}
\fancyhf{}
\lhead{\textsc{PCSI -- Analyse}}
\chead{\textsc{DS : Suites, continuité, dérivation}}
\rhead{Page \thepage/\pageref{LastPage}}
\cfoot{}

\titleformat{\section}
{\Large\bfseries\color{mainblue}}
{\thesection}{0.8em}{}

\titleformat{\subsection}
{\large\bfseries\color{darkred}}
{\thesubsection}{0.8em}{}

\titleformat{\subsubsection}
{\normalsize\bfseries\color{darkgreen}}
{\thesubsubsection}{0.8em}{}

% ============================================================
% ENVIRONNEMENTS
% ============================================================

\newtcolorbox{consignesbox}{
    colback=lightblue,
    colframe=mainblue,
    title=\textbf{Consignes générales},
    fonttitle=\bfseries,
    arc=2mm,
    boxrule=0.8pt
}

\newtcolorbox{notionbox}{
    colback=lightgreen,
    colframe=darkgreen,
    title=\textbf{Notion introduite dans le sujet},
    fonttitle=\bfseries,
    arc=2mm,
    boxrule=0.8pt
}

\newtcolorbox{attentionbox}{
    colback=lightred,
    colframe=darkred,
    title=\textbf{Attention},
    fonttitle=\bfseries,
    arc=2mm,
    boxrule=0.8pt
}

\newtcolorbox{baremebox}{
    colback=lightgray,
    colframe=black!60,
    title=\textbf{Barème indicatif},
    fonttitle=\bfseries,
    arc=2mm,
    boxrule=0.8pt
}

\newcommand{\R}{\mathbb{R}}
\newcommand{\N}{\mathbb{N}}
\newcommand{\eps}{\varepsilon}
\newcommand{\dd}{\,\mathrm{d}}

\setlist[itemize]{label=\(\triangleright\), leftmargin=1.2cm}
\setlist[enumerate]{leftmargin=1cm}

% ============================================================
% DOCUMENT
% ============================================================

\begin{document}

% ============================================================
% PAGE DE TITRE
% ============================================================

\begin{center}
    {\Huge \bfseries Devoir surveillé d'analyse}\\[0.3cm]
    {\Large \bfseries Suites, continuité, dérivation}\\[0.2cm]
    {\large Niveau PCSI -- Esprit CCP-INP faible}\\[0.4cm]
    \rule{0.85\linewidth}{0.5pt}\\[0.3cm]
    \textbf{Durée conseillée : 4 heures}\\
    \textbf{Calculatrice interdite. Documents interdits.}\\[0.2cm]
    \rule{0.85\linewidth}{0.5pt}
\end{center}

\vspace{0.5cm}

\begin{consignesbox}
La qualité de la rédaction, la précision des quantificateurs, la vérification des hypothèses des théorèmes utilisés et la clarté des conclusions seront fortement prises en compte.

On attachera une grande importance aux points suivants :
\begin{itemize}
    \item annoncer clairement les théorèmes utilisés ;
    \item vérifier leurs hypothèses ;
    \item ne pas confondre existence, unicité et calcul effectif ;
    \item justifier les passages à la limite ;
    \item distinguer les arguments de monotonie, de continuité, de dérivabilité et de majoration.
\end{itemize}

Les deux problèmes principaux sont largement indépendants.  
Certaines questions demandent volontairement de choisir soi-même la méthode : aucune indication directe n'est alors donnée, afin d'évaluer la capacité à reconnaître un raisonnement du cours.
\end{consignesbox}

\vspace{0.3cm}

\begin{attentionbox}
Aucune notion hors programme n'est nécessaire.  
En particulier, il est interdit d'utiliser :
\[
\text{intégrales, séries, développements limités, formule de Taylor, convexité, équations différentielles.}
\]
La seule approximation locale autorisée est celle qui découle directement de la définition de la dérivabilité :
\[
f(a+h)=f(a)+hf'(a)+h\eps(h),
\qquad \eps(h)\xrightarrow[h\to 0]{}0.
\]
\end{attentionbox}

\newpage

% ============================================================
% PROBLEME I
% ============================================================

\section*{Problème I -- La racine de \(2\) comme objet d'ordre, puis comme limite d'un algorithme}

Ce problème étudie la construction de la racine carrée de \(2\) à partir de la borne supérieure, puis un algorithme permettant de l'approcher.

On pose
\[
A=\{x\in \R_+ \mid x^2\leq 2\}.
\]

On admet uniquement les propriétés usuelles de l'ordre sur \(\R\), la propriété de la borne supérieure, les résultats du cours sur les suites monotones et la continuité des fonctions usuelles.

% ------------------------------------------------------------

\subsection*{Partie A -- Construction par borne supérieure}

\begin{enumerate}
    \item Montrer que \(A\) est non vide et majoré.

    \item On note
    \[
    r=\sup A.
    \]
    Montrer que
    \[
    1\leq r\leq 2.
    \]

    \item Montrer que l'on ne peut pas avoir \(r^2<2\).

    \medskip

    \emph{Indication volontairement minimale : raisonner à partir de la définition de la borne supérieure et construire un réel strictement plus grand que \(r\) appartenant encore à \(A\).}

    \item Montrer que l'on ne peut pas avoir \(r^2>2\).

    \medskip

    \emph{Indication volontairement minimale : utiliser l'idée symétrique, mais attention : il faut montrer qu'un nombre strictement inférieur à \(r\) est encore un majorant de \(A\).}

    \item Conclure que
    \[
    r^2=2.
    \]

    \item Montrer que \(r\) est l'unique réel positif vérifiant \(x^2=2\).
\end{enumerate}

Dans toute la suite du problème, on note
\[
\sqrt{2}=r.
\]

% ------------------------------------------------------------

\subsection*{Partie B -- Une fonction itérative}

On définit la fonction
\[
\varphi : ]0,+\infty[ \longrightarrow ]0,+\infty[
\]
par
\[
\varphi(x)=\frac12\left(x+\frac{2}{x}\right).
\]

\begin{enumerate}[resume]
    \item Montrer que \(\varphi\) est bien définie et continue sur \(]0,+\infty[\).

    \item Montrer que, pour tout \(x>0\),
    \[
    \varphi(x)-\sqrt{2}
    =
    \frac{(x-\sqrt{2})^2}{2x}.
    \]

    \item En déduire que
    \[
    \forall x>0,\qquad \varphi(x)\geq \sqrt{2}.
    \]

    \item Montrer que, pour tout \(x\geq \sqrt{2}\),
    \[
    \varphi(x)\leq x.
    \]

    \item En déduire que l'intervalle \([\sqrt{2},+\infty[\) est stable par \(\varphi\).
\end{enumerate}

% ------------------------------------------------------------

\subsection*{Partie C -- Convergence de la suite de Héron}

\begin{notionbox}
La suite étudiée ci-dessous est appelée \emph{suite de Héron}. Elle fournit un algorithme d'approximation de \(\sqrt{2}\). Aucune connaissance préalable sur cet algorithme n'est attendue.
\end{notionbox}

On définit la suite \((u_n)_{n\in\N}\) par
\[
u_0=2,
\qquad
u_{n+1}=\varphi(u_n).
\]

\begin{enumerate}[resume]
    \item Montrer que, pour tout \(n\in\N\),
    \[
    u_n\geq \sqrt{2}.
    \]

    \item Montrer que la suite \((u_n)\) est décroissante.

    \item En déduire que \((u_n)\) converge vers un réel \(\ell\).

    \item Montrer que
    \[
    \ell=\sqrt{2}.
    \]
    On prendra soin de justifier le passage à la limite dans la relation de récurrence.

    \item Montrer que, pour tout \(n\in\N\),
    \[
    0\leq u_{n+1}-\sqrt{2}
    \leq \frac12\left(u_n-\sqrt{2}\right)^2.
    \]

    \item Sans chercher à donner une formule explicite de \(u_n\), expliquer pourquoi cette inégalité rend plausible la rapidité de convergence de l'algorithme.
\end{enumerate}

% ------------------------------------------------------------

\subsection*{Partie D -- Deux suites adjacentes cachées}

On pose, pour tout \(n\in\N\),
\[
v_n=\frac{2}{u_n}.
\]

\begin{enumerate}[resume]
    \item Montrer que, pour tout \(n\in\N\),
    \[
    0<v_n\leq \sqrt{2}\leq u_n.
    \]

    \item Montrer que \((v_n)\) est croissante.

    \item Montrer que
    \[
    u_n-v_n \longrightarrow 0.
    \]

    \item En déduire que les suites \((u_n)\) et \((v_n)\) ont la même limite.

    \item Retrouver ainsi la limite commune de \((u_n)\) et \((v_n)\).
\end{enumerate}

\newpage

% ============================================================
% PROBLEME II
% ============================================================

\section*{Problème II -- Une fonction prolongée, une bijection et des inégalités par accroissements finis}

Dans ce problème, on étudie une fonction définie d'abord par un quotient, puis prolongée par continuité. On montrera ensuite qu'elle définit une bijection et que sa réciproque peut être contrôlée par le théorème des accroissements finis.

On considère la fonction \(f\) définie sur \(]-1,+\infty[\) par
\[
f(x)=
\begin{cases}
\dfrac{\sqrt{1+x}-1}{x} & \text{si } x\neq 0,\\[1.2em]
\dfrac12 & \text{si } x=0.
\end{cases}
\]

% ------------------------------------------------------------

\subsection*{Partie A -- Prolongement et premières propriétés}

\begin{enumerate}
    \item Montrer que, pour tout \(x\in]-1,+\infty[\setminus\{0\}\),
    \[
    f(x)=\frac{1}{\sqrt{1+x}+1}.
    \]

    \item En déduire que \(f\) est continue en \(0\).

    \item Montrer que \(f\) est continue sur \(]-1,+\infty[\).

    \item Déterminer les limites suivantes :
    \[
    \lim_{x\to -1^+} f(x),
    \qquad
    \lim_{x\to +\infty} f(x).
    \]

    \item Montrer que
    \[
    0<f(x)<1
    \]
    pour tout \(x\in]-1,+\infty[\).
\end{enumerate}

% ------------------------------------------------------------

\subsection*{Partie B -- Dérivabilité et variations}

\begin{enumerate}[resume]
    \item Montrer que \(f\) est dérivable sur \(]-1,+\infty[\), y compris en \(0\), et calculer \(f'(x)\).

    \item Montrer que \(f\) est strictement décroissante sur \(]-1,+\infty[\).

    \item Retrouver la stricte décroissance de \(f\) sans utiliser directement le signe de \(f'\), mais en utilisant le théorème de Rolle.

    \medskip

    \emph{Cette question vise à faire réapparaître une idée de démonstration du cours.}

    \item Montrer que \(f\) réalise une bijection de \(]-1,+\infty[\) sur \(]0,1[\).

    \item Déterminer explicitement la fonction réciproque \(f^{-1}\).
\end{enumerate}

% ------------------------------------------------------------

\subsection*{Partie C -- Contrôle de la réciproque}

On note
\[
g=f^{-1}.
\]

\begin{enumerate}[resume]
    \item Donner le domaine de définition de \(g\) et son expression explicite.

    \item Montrer que \(g\) est dérivable sur \(]0,1[\), puis calculer \(g'(y)\).

    \item Soient \(0<a<b<1\). Montrer qu'il existe une constante \(M_{a,b}>0\) telle que
    \[
    \forall y,z\in[a,b],
    \qquad
    |g(y)-g(z)|\leq M_{a,b}|y-z|.
    \]

    \item Donner une valeur possible de \(M_{a,b}\).

    \item Expliquer pourquoi il est impossible de choisir une constante \(M>0\) qui convienne sur tout l'intervalle \(]0,1[\).

    \medskip

    \emph{On attend ici un raisonnement mathématique, pas seulement une intuition graphique.}
\end{enumerate}

% ------------------------------------------------------------

\subsection*{Partie D -- Une inégalité fine sans développement limité}

\begin{enumerate}[resume]
    \item Montrer que, pour tout \(x>-1\),
    \[
    f'(x)=-\frac{1}{2\sqrt{1+x}(\sqrt{1+x}+1)^2}.
    \]

    \item Soit \(x\geq 0\). En appliquant l'inégalité des accroissements finis à \(f\) sur \([0,x]\), montrer que
    \[
    0\leq \frac12-f(x)\leq \frac{x}{8}.
    \]

    \item Traduire cette inégalité sous la forme
    \[
    0\leq \frac12-\frac{\sqrt{1+x}-1}{x}\leq \frac{x}{8}
    \qquad (x>0).
    \]

    \item En déduire que, pour tout \(x>0\),
    \[
    0\leq 1+\frac{x}{2}-\sqrt{1+x}\leq \frac{x^2}{8}.
    \]

    \item Montrer directement, à partir de la définition de la dérivabilité de la fonction \(x\mapsto \sqrt{1+x}\) en \(0\), que
    \[
    \sqrt{1+x}=1+\frac{x}{2}+x\eps(x),
    \]
    où
    \[
    \eps(x)\longrightarrow 0
    \qquad \text{quand } x\to 0.
    \]
    Comparer ce résultat avec celui obtenu à la question précédente.
\end{enumerate}

\newpage

% ============================================================
% PROBLEME III BONUS
% ============================================================

% ============================================================
% PROBLEME III
% ============================================================

\section*{Problème III -- Dichotomie, théorème des valeurs intermédiaires et approximation certifiée d'une racine}

Ce problème prolonge le théorème des valeurs intermédiaires sous une forme constructive.  
L'objectif est de comprendre comment une preuve d'existence peut devenir un algorithme d'approximation rigoureux.

\begin{notionbox}
On appelle \emph{procédure de dichotomie} la construction suivante.

Soit \(h\) une fonction continue sur un segment \([a,b]\), telle que
\[
h(a)\leq 0\leq h(b).
\]
On construit deux suites \((a_n)\) et \((b_n)\) ainsi :
\[
a_0=a,
\qquad
b_0=b.
\]
Une fois \(a_n\) et \(b_n\) construits, on pose
\[
m_n=\frac{a_n+b_n}{2}.
\]
Si
\[
h(m_n)\leq 0,
\]
on pose
\[
a_{n+1}=m_n,
\qquad
b_{n+1}=b_n.
\]
Sinon, on pose
\[
a_{n+1}=a_n,
\qquad
b_{n+1}=m_n.
\]

On dit que \([a_n,b_n]\) est le \emph{\(n\)-ième intervalle de localisation}.
\end{notionbox}

% ------------------------------------------------------------

\subsection*{Partie A -- Une preuve constructive du théorème des valeurs intermédiaires}

Dans cette partie, \(h\) est seulement supposée continue sur \([a,b]\), avec
\[
h(a)\leq 0\leq h(b).
\]

\begin{enumerate}
    \item Montrer que, pour tout \(n\in\N\),
    \[
    a_n\leq b_n.
    \]

    \item Montrer que, pour tout \(n\in\N\),
    \[
    h(a_n)\leq 0\leq h(b_n).
    \]

    \item Montrer que la suite \((a_n)\) est croissante et que la suite \((b_n)\) est décroissante.

    \item Montrer que, pour tout \(n\in\N\),
    \[
    b_{n+1}-a_{n+1}
    =
    \frac12(b_n-a_n).
    \]

    \item En déduire que, pour tout \(n\in\N\),
    \[
    b_n-a_n=\frac{b-a}{2^n}.
    \]

    \item Montrer que les suites \((a_n)\) et \((b_n)\) sont adjacentes.

    \item On note \(\alpha\) leur limite commune. Montrer que
    \[
    \alpha\in [a,b].
    \]

    \item Montrer que
    \[
    h(\alpha)=0.
    \]

    \medskip

    \emph{Cette question doit être rédigée avec soin : on attend un vrai passage à la limite utilisant la continuité de \(h\).}

    \item Expliquer en quelques lignes pourquoi cette procédure donne une démonstration constructive du théorème des valeurs intermédiaires dans le cas où \(h(a)\leq 0\leq h(b)\).
\end{enumerate}

% ------------------------------------------------------------

\subsection*{Partie B -- Cas d'une racine unique : encadrement et erreur}

On suppose désormais que \(h\) est continue sur \([a,b]\), que
\[
h(a)\leq 0\leq h(b),
\]
et que l'équation
\[
h(x)=0
\]
admet une unique solution dans \([a,b]\). On note cette solution \(\alpha\).

\begin{enumerate}[resume]
    \item Montrer que, pour tout \(n\in\N\),
    \[
    a_n\leq \alpha\leq b_n.
    \]

    \item On définit le milieu de l'intervalle de localisation par
    \[
    c_n=\frac{a_n+b_n}{2}.
    \]
    Montrer que, pour tout \(n\in\N\),
    \[
    |c_n-\alpha|\leq \frac{b_n-a_n}{2}.
    \]

    \item En déduire que
    \[
    |c_n-\alpha|\leq \frac{b-a}{2^{n+1}}.
    \]

    \item Montrer que la suite \((c_n)\) converge vers \(\alpha\).

    \item Soit \(\varepsilon>0\). Montrer qu'il existe un rang \(N\in\N\) tel que
    \[
    \forall n\geq N,
    \qquad
    |c_n-\alpha|\leq \varepsilon.
    \]

    \medskip

    \emph{On ne demande pas nécessairement d'exprimer \(N\) à l'aide d'un logarithme. Une justification par la divergence de \(2^n\) suffit.}

    \item Interpréter le résultat précédent : en quoi la dichotomie fournit-elle une approximation \emph{certifiée} de la racine ?
\end{enumerate}

% ------------------------------------------------------------

\subsection*{Partie C -- Application : la racine plastique}

On considère la fonction
\[
p : [1,2]\longrightarrow \R,
\qquad
p(x)=x^3-x-1.
\]

La racine réelle positive de l'équation
\[
x^3=x+1
\]
est parfois appelée \emph{racine plastique}. Aucune connaissance sur ce nombre n'est attendue.

\begin{enumerate}[resume]
    \item Montrer que \(p\) est continue sur \([1,2]\).

    \item Calculer \(p(1)\) et \(p(2)\). En déduire que l'équation
    \[
    p(x)=0
    \]
    admet au moins une solution dans \([1,2]\).

    \item Montrer que \(p\) est strictement croissante sur \([1,2]\).

    \item En déduire que l'équation
    \[
    x^3-x-1=0
    \]
    admet une unique solution dans \([1,2]\). On la note \(\rho\).

    \item On applique la procédure de dichotomie à \(p\) sur \([1,2]\). Déterminer explicitement les trois premiers intervalles de localisation :
    \[
    [a_0,b_0],\qquad [a_1,b_1],\qquad [a_2,b_2],\qquad [a_3,b_3].
    \]

    \item Vérifier que l'on obtient
    \[
    \frac54\leq \rho\leq \frac{11}{8}.
    \]

    \item Poursuivre une étape supplémentaire et montrer que
    \[
    \frac{21}{16}\leq \rho\leq \frac{11}{8}.
    \]

    \item Donner un encadrement de \(\rho\) par un intervalle de longueur inférieure ou égale à
    \[
    \frac{1}{16}.
    \]

    \item Sans calculer davantage de valeurs de \(p\), déterminer un nombre d'étapes suffisant pour garantir une approximation de \(\rho\) à \(10^{-3}\) près par le milieu de l'intervalle.

    \medskip

    \emph{On pourra utiliser le fait que}
    \[
    2^{10}=1024.
    \]

    \item On note \(c_n\) le milieu du \(n\)-ième intervalle de localisation. Montrer que
    \[
    c_n\longrightarrow \rho.
    \]
\end{enumerate}

% ------------------------------------------------------------

\subsection*{Partie D -- Une question de réflexion : existence, unicité et algorithme}

\begin{enumerate}[resume]
    \item Dans la partie A, l'unicité de la racine était-elle nécessaire pour construire les suites \((a_n)\) et \((b_n)\) ?

    \item Dans la partie B, expliquer précisément où l'hypothèse d'unicité est utilisée.

    \item Donner un exemple simple de fonction continue \(h\) sur un segment \([a,b]\), vérifiant
    \[
    h(a)\leq 0\leq h(b),
    \]
    mais possédant plusieurs zéros dans \([a,b]\).

    \item Dans ce cas, la procédure de dichotomie converge-t-elle forcément vers un zéro de \(h\) ? Justifier la réponse.

    \item Expliquer pourquoi la dichotomie est plus robuste qu'une méthode fondée uniquement sur le calcul formel exact d'une racine.
\end{enumerate}

\newpage

% ============================================================
% BAREME
% ============================================================

\section*{Barème indicatif sur 100 points}

\begin{baremebox}
La note finale n'est qu'un indicateur. Le barème valorise surtout :
\begin{itemize}
    \item le choix d'une méthode adaptée ;
    \item la vérification des hypothèses ;
    \item la rigueur des quantificateurs ;
    \item la rédaction propre ;
    \item la justification des limites ;
    \item le bon usage des théorèmes ;
    \item la conclusion clairement formulée.
\end{itemize}
\end{baremebox}

% ------------------------------------------------------------

% ------------------------------------------------------------

\subsection*{Problème I -- 35 points}

\subsubsection*{Partie A -- 10 points}

\begin{tabularx}{\linewidth}{>{\bfseries}c X c}
\toprule
Question & Attendu principal & Points \\
\midrule
1 & \(0\in A\), \(A\) majoré, par exemple par \(2\). & 1,5 \\
2 & \(1\in A\), donc \(r\geq 1\), et \(2\) est majorant. & 1,5 \\
3 & Construire \(\eta>0\) tel que \((r+\eta)^2<2\), contradiction. & 2 \\
4 & Construire \(\eta>0\) tel que \(r-\eta\) soit encore majorant, contradiction. & 2,5 \\
5 & Conclure \(r^2=2\). & 1 \\
6 & Unicité par stricte croissance de \(x\mapsto x^2\) sur \(\R_+\). & 1,5 \\
\bottomrule
\end{tabularx}

\subsubsection*{Partie B -- 8 points}

\begin{tabularx}{\linewidth}{>{\bfseries}c X c}
\toprule
Question & Attendu principal & Points \\
\midrule
7 & Définition et continuité de \(\varphi\). & 1 \\
8 & Identité fondamentale. & 2 \\
9 & Minoration par \(\sqrt2\). & 1 \\
10 & Inégalité \(\varphi(x)\leq x\) pour \(x\geq\sqrt2\). & 2 \\
11 & Stabilité de \([\sqrt2,+\infty[\). & 2 \\
\bottomrule
\end{tabularx}

\subsubsection*{Partie C -- 10 points}

\begin{tabularx}{\linewidth}{>{\bfseries}c X c}
\toprule
Question & Attendu principal & Points \\
\midrule
12 & Récurrence et stabilité. & 1,5 \\
13 & Décroissance. & 1,5 \\
14 & Convergence par suite monotone bornée. & 1,5 \\
15 & Passage à la limite, point fixe, unicité. & 2,5 \\
16 & Inégalité d'erreur. & 2 \\
17 & Interprétation qualitative. & 1 \\
\bottomrule
\end{tabularx}

\subsubsection*{Partie D -- 7 points}

\begin{tabularx}{\linewidth}{>{\bfseries}c X c}
\toprule
Question & Attendu principal & Points \\
\midrule
18 & Encadrement de \(v_n\). & 1,5 \\
19 & Croissance de \((v_n)\). & 1,5 \\
20 & Limite de \(u_n-v_n\). & 1,5 \\
21 & Adjacence. & 1,5 \\
22 & Limite commune. & 1 \\
\bottomrule
\end{tabularx}

% ------------------------------------------------------------

\subsection*{Problème II -- 40 points}

\subsubsection*{Partie A -- 8 points}

\begin{tabularx}{\linewidth}{>{\bfseries}c X c}
\toprule
Question & Attendu principal & Points \\
\midrule
1 & Rationalisation correcte. & 1,5 \\
2 & Continuité en \(0\). & 1,5 \\
3 & Continuité globale. & 1 \\
4 & Limites aux bornes. & 2 \\
5 & Encadrement \(0<f(x)<1\). & 2 \\
\bottomrule
\end{tabularx}

\subsubsection*{Partie B -- 11 points}

\begin{tabularx}{\linewidth}{>{\bfseries}c X c}
\toprule
Question & Attendu principal & Points \\
\midrule
6 & Dérivabilité et calcul de \(f'\). & 3 \\
7 & Stricte décroissance. & 2 \\
8 & Argument par Rolle. & 2 \\
9 & Bijection continue monotone. & 2,5 \\
10 & Expression explicite de \(f^{-1}\). & 1,5 \\
\bottomrule
\end{tabularx}

\subsubsection*{Partie C -- 9 points}

\begin{tabularx}{\linewidth}{>{\bfseries}c X c}
\toprule
Question & Attendu principal & Points \\
\midrule
11 & Domaine et expression de \(g\). & 1,5 \\
12 & Dérivabilité et calcul de \(g'\). & 2 \\
13 & Contrôle sur \([a,b]\) par TAF ou IAF. & 2 \\
14 & Constante possible \(M_{a,b}\). & 1,5 \\
15 & Impossibilité globale car \(g'\) explose vers \(0^+\). & 2 \\
\bottomrule
\end{tabularx}

\subsubsection*{Partie D -- 12 points}

\begin{tabularx}{\linewidth}{>{\bfseries}c X c}
\toprule
Question & Attendu principal & Points \\
\midrule
16 & Expression de \(f'\). & 2 \\
17 & Application correcte de l'IAF. & 3 \\
18 & Reformulation de l'inégalité. & 1,5 \\
19 & Approximation contrôlée de \(\sqrt{1+x}\). & 2,5 \\
20 & Dérivabilité en \(0\) par définition et comparaison. & 3 \\
\bottomrule
\end{tabularx}

% ------------------------------------------------------------

\subsection*{Problème III -- 25 points}

\subsubsection*{Partie A -- 8 points}

\begin{tabularx}{\linewidth}{>{\bfseries}c X c}
\toprule
Question & Attendu principal & Points \\
\midrule
1 & Montrer \(a_n\leq b_n\) par récurrence. & 0,75 \\
2 & Conservation de l'encadrement de signe. & 1,25 \\
3 & Monotonie de \((a_n)\) et \((b_n)\). & 1 \\
4 & Relation sur la longueur des intervalles. & 1 \\
5 & Formule \(b_n-a_n=(b-a)/2^n\). & 1 \\
6 & Adjacence des deux suites. & 1 \\
7 & Localisation de la limite commune. & 0,75 \\
8 & Passage à la limite et continuité de \(h\). & 1 \\
9 & Interprétation constructive du TVI. & 0,25 \\
\bottomrule
\end{tabularx}

\subsubsection*{Partie B -- 5 points}

\begin{tabularx}{\linewidth}{>{\bfseries}c X c}
\toprule
Question & Attendu principal & Points \\
\midrule
10 & Encadrement \(a_n\leq\alpha\leq b_n\). & 1 \\
11 & Estimation \(|c_n-\alpha|\leq (b_n-a_n)/2\). & 1 \\
12 & Estimation explicite avec \(2^{n+1}\). & 1 \\
13 & Convergence de \((c_n)\) vers \(\alpha\). & 1 \\
14 & Traduction en précision \(\varepsilon\). & 0,75 \\
15 & Interprétation d'une approximation certifiée. & 0,25 \\
\bottomrule
\end{tabularx}

\subsubsection*{Partie C -- 8 points}

\begin{tabularx}{\linewidth}{>{\bfseries}c X c}
\toprule
Question & Attendu principal & Points \\
\midrule
16 & Continuité de \(p\). & 0,5 \\
17 & Signes en \(1\) et \(2\), existence par TVI. & 1 \\
18 & Stricte croissance de \(p\) sur \([1,2]\). & 1 \\
19 & Unicité de la racine. & 1 \\
20 & Trois premiers intervalles de dichotomie. & 1,5 \\
21 & Encadrement \(\frac54\leq\rho\leq\frac{11}{8}\). & 0,75 \\
22 & Étape supplémentaire \(\frac{21}{16}\leq\rho\leq\frac{11}{8}\). & 0,75 \\
23 & Intervalle de longueur \(\leq 1/16\). & 0,5 \\
24 & Nombre d'étapes pour une précision \(10^{-3}\). & 0,75 \\
25 & Convergence des milieux \(c_n\) vers \(\rho\). & 0,25 \\
\bottomrule
\end{tabularx}

\subsubsection*{Partie D -- 4 points}

\begin{tabularx}{\linewidth}{>{\bfseries}c X c}
\toprule
Question & Attendu principal & Points \\
\midrule
26 & L'unicité n'est pas nécessaire à la construction. & 0,75 \\
27 & Identifier où l'unicité est utilisée. & 0,75 \\
28 & Exemple d'une fonction avec plusieurs zéros. & 0,75 \\
29 & Justifier que la limite obtenue est encore un zéro. & 1 \\
30 & Discussion sur la robustesse de la dichotomie. & 0,75 \\
\bottomrule
\end{tabularx}

\newpage




\end{document}